\section{Thiết kế Edge Runtime}

AI Runtime là một container độc lập do Agent tạo ra. Toàn bộ cấu hình được
truyền qua biến môi trường; khóa \texttt{RUNTIME\_CONFIG} (JSON) chọn nguồn
khung hình (\texttt{source}), FPS mục tiêu (\texttt{target\_fps}, mặc định 15),
detector (\texttt{ultralytics} khi có \texttt{MODEL\_PATH}, ngược lại
\texttt{synthetic}), ngưỡng tin cậy (\texttt{conf}) và chu kỳ báo cáo. Luồng xử
lý được thể hiện trong Hình~\ref{fig:edge-runtime-design}.

\begin{figure}[H]
\centering
\resizebox{\textwidth}{!}{%
\begin{tikzpicture}[node distance=0.55cm and 0.7cm]
  \node[flowbox] (camera) {Camera\\RTSP/USB/file};
  \node[flowbox, right=of camera] (capture) {Capture task\\(FrameSource)};
  \node[flowbox, right=of capture, fill=orange!8] (queue) {Hàng đợi\\8 khung};
  \node[flowbox, right=of queue] (infer) {Inference task\\(Detector)};
  \node[flowbox, right=of infer, text width=3cm] (metrics) {Inference\\Metrics};
  \node[flowbox, below=1.4cm of metrics, text width=3cm] (report) {Report task\\(5 s)};
  \node[flowbox, left=of report, fill=green!6] (agent) {Edge Agent\\HTTP API};
  \draw[flowarrow] (camera) -- (capture);
  \draw[flowarrow] (capture) -- (queue);
  \draw[flowarrow] (queue) -- (infer);
  \draw[flowarrow] (infer) -- (metrics);
  \draw[flowarrow] (metrics) -- (report);
  \draw[flowarrow] (report) -- node[font=\scriptsize, above]{telemetry} (agent);
  \draw[flowarrow, dashed] (capture.south) -- ++(0,-0.6) -| node[pos=0.25, above, font=\scriptsize]{frame drop, sự kiện camera} ([xshift=-0.6cm]metrics.south);
\end{tikzpicture}}
\caption{Thiết kế pipeline Edge Runtime}
\label{fig:edge-runtime-design}
\end{figure}

\textbf{Ba tác vụ song song.} InferenceService chạy ba tác vụ trong một
\texttt{TaskGroup}: \emph{capture} đọc khung từ FrameSource; \emph{infer} lấy khung
từ hàng đợi và gọi Detector với thời gian chờ tối đa 2~s; \emph{report} định kỳ
chụp ảnh chỉ số và gửi cho Agent. Khi một tác vụ lỗi nghiêm trọng, cả nhóm dừng
và container thoát, để vòng reconcile của Agent tạo lại container.

\textbf{Hàng đợi giới hạn, ưu tiên khung mới.} Hàng đợi chỉ giữ 8 khung. Khi
đầy, khung cũ nhất bị loại và được tính là \emph{frame drop}. Chính sách này giữ
độ trễ không tăng vô hạn khi detector chậm hơn camera: hệ thống chấp nhận bỏ
khung để luôn suy luận trên dữ liệu mới nhất.

\textbf{Nguồn khung và detector là port.} \texttt{OpenCvFrameSource} đọc RTSP,
USB (chỉ số thiết bị) hoặc tệp video, tự kết nối lại sau 3~s khi luồng mất và
phát sự kiện \texttt{camera.disconnected}. \texttt{UltralyticsDetector} nạp mô
hình qua API \texttt{YOLO(...)} của Ultralytics nên dùng được trực tiếp tệp
\texttt{.pt}, \texttt{.onnx} hoặc TensorRT \texttt{.engine}; lời gọi suy luận
chạy trong thread riêng để không chặn vòng sự kiện. Các adapter
\texttt{SyntheticFrameSource} và \texttt{SyntheticDetector} cho phép chạy toàn bộ
pipeline khi không có camera hoặc GPU.

\textbf{Chỉ số runtime.} \texttt{InferenceMetrics} tích lũy trong mỗi cửa sổ báo
cáo: FPS suy luận và FPS camera, latency trung bình và p50/p95/p99, tỷ lệ lỗi và
timeout, tỷ lệ frame drop, tuổi khung mới nhất, thời gian camera bị đứng hình,
độ sáng trung bình, số detection, confidence trung bình/độ lệch/nhỏ nhất và
entropy phân bố lớp. Các trường này khớp tên với đặc tả telemetry ở
Mục~\ref{sec:monitoring-design}.
